\noindent\textbf{Elección del procedimiento.} Comprobar primero Reynolds. Si el modelo laminar es apropiado, el factor de Darcy es $64/\mathrm{Re}$; introducirlo en Darcy--Weisbach da la pérdida de presión. La potencia es presión por caudal.

\begin{enumerate}

\item Usar velocidad media y diámetro interior.

\[\fitmath{\mathrm{Re}=\frac{\rho\cdot V\cdot D}{\mu}}\]

\[\fitmath{\mathrm{Re}=\frac{(870)\cdot (1.2)\cdot (0.004)}{(0.012)}\approx 348\,}\]

\item Este es el factor de Darcy; el de Fanning sería cuatro veces menor.

\[\fitmath{\mathrm{fD}=\frac{64}{\mathrm{Re}}}\]

\[\fitmath{\mathrm{fD}=\frac{64}{(348)}\approx 0.183908\,}\]

\item La longitud se divide por el diámetro, no por el radio.

\[\fitmath{\mathrm{dp}=\frac{\frac{\mathrm{fD}\cdot L}{D}\cdot \rho\cdot {V}^{2}}{2}}\]

\[\fitmath{\mathrm{dp}=\frac{\frac{(0.183908046)\cdot (2)}{(0.004)}\cdot (870)\cdot {(1.2)}^{2}}{2}\approx 57600\,\mathrm{Pa}}\]

\item Obtener el caudal desde la velocidad media.

\[\fitmath{Q=\frac{\pi\cdot {D}^{2}\cdot V}{4}}\]

\[\fitmath{Q=\frac{\pi\cdot {(0.004)}^{2}\cdot (1.2)}{4}\approx 1.507964\times10^{-5}\,\mathrm{m^{3}/s}}\]

\item Convertir la pérdida de energía por volumen en pérdida por segundo.

\[\fitmath{P=\mathrm{dp}\cdot Q}\]

\[\fitmath{P=(57600)\cdot (1.507964474\times10^{-5})\approx 0.8685875\,\mathrm{W}}\]

\end{enumerate}
